\documentclass[11pt,a4paper]{article}
\usepackage[margin=25mm]{geometry}
\usepackage{amsmath,amssymb,amsthm,booktabs,longtable,hyperref,xr}
\externaldocument[U-]{umbrella}
\newtheorem{proposition}{Proposition}
\newtheorem{corollary}[proposition]{Corollary}
\newcommand{\Sym}{\operatorname{Sym}}
\newcommand{\Aut}{\operatorname{Aut}}
\title{Shadow intersections in sporadic and tree geometries}
\author{A complete corrected principal-result assembly}
\date{}
\begin{document}\maketitle
\begin{abstract}
One residue theorem and one complete witness search determine the
geometric intersection properties of the supplied sporadic geometries.
We consolidate the previously reviewed finite calculations, their actual
geometry hypotheses and their exact family correspondences. Of the
57 numerical rows printed by Jacobs and Leemans, 53 agree with the
current supplied-roster or retained-census calculations. The four
differences concern two repeated rank-three reports: $M_{12}$ and $J_3$.
We also recover the Suzuki shadow split and the uniform tree criterion.
A complete companion proves the available historical name correspondences.
For $J_1$ we state the complete mathematical classification by explicit
type sizes, reserving three unverified historical ordinal correspondences.
This is a corrected principal-result recovery, not literal table-name
equivalence.
\end{abstract}

\section{A single finite proof boundary}
Our source is the complete sixteen-page edition of Jacobs and Leemans
\cite{jl}, together with its twenty-file supplementary geometry archive.
The common paper \emph{Witnesses and orbits in finite incidence geometries}
is a whole mathematical dependency. Its covering, orbit, double-coset,
residue, sparse-coset, maximality and transporter arguments are written
once there. The original orbit method is attributed to Jacobs and
Leemans; the finite-group and block arguments are standard. Our common
implementation and this assembly do not claim new priority for the
source results.

For an enlarged Topic D, its entire accepted common paper
\emph{Stabilizers, extensions and the enumeration of finite geometries}
remains a paid dependency. Its coset realization, panel recognition,
local primitivity and augmentation theorems already contain part of
the infrastructure used here. The additional shadow theory concerns
geometric witness flags, normalized object--flag queries, sparse coset
representatives and the rank-by-rank existential quantifier; these are
not supplied by the old support-counting theorem. The detailed overlap
ledger also marks repeated standard proofs as overlap. No existing
common section is removed, and no overlap discount or compression
factor is inferred from this assembly.

An input is an ordered family of actual subgroups in a specified acting
group. The equivalence used for a finite roster is one automorphism of
that group applied to the \emph{whole} family, followed by a permutation
of types. Equality of subgroup orders, intersection profiles, shadow
answers, or aggregate counts never establishes this equivalence.
The last active assignment in each archive file is parsed literally;
it is kept distinct from the order of construction and from a historical
geometry label.

\begin{proposition}\label{application-method}
Suppose a finite roster is given explicitly, and each input either passes
the common geometry tests or is identified by an actual whole-family
isomorphism with a previously accepted marked input having those
properties. For every eligible equivalence-class representative, test
all base residues of rank at least two and every remaining type by the
common shadow theorem. The resulting all/exists quantifiers give exactly
its $(\mathrm{IP})_k$, $(\mathrm{WIP})_k$, IP and WIP values. Summing
these values over the exact family quotient gives the corresponding
roster counts. A separate prior classification is needed to assert
that the roster contains every geometry of the group.
\end{proposition}
\begin{proof}
The eligible-extension equalities establish full flag transitivity, not
merely transitivity on pairs. The common residue theorem then identifies
the residue at each base flag with its actual intersection-subgroup
coset geometry, and transports it to every flag of the same type. In
each residue the orbit/double-coset theorem exhausts all normalized
element--flag queries and all eligible witness flags. Its sparse
implementation changes only the representation of right cosets; empty
flags still use the exact whole-universe cardinality. Thus each recorded
type decision is the required equivalence, including negative decisions.
Taking all types gives IP at that residue; taking at least one gives
WIP, with the successful type allowed to vary between residues.
Taking all residues of rank $k$, and then all $k\geq2$, gives the stated
properties. Panels are separately checked for geometry hypotheses;
rank zero is never used for the existential WIP condition.

An actual family automorphism and type permutation induce an incidence
isomorphism, which transports flags, shadows and witness flags. Therefore
every tested property is constant on a family class. The exact
transporter criterion makes the quotient neither omit a class nor count
one twice. Summing its class values proves the roster counts, but gives
no assertion about an input absent from the roster. This is precisely
why the completeness premise is separate.
\end{proof}

The source population requires flag transitivity, firmness, residual
connectedness, residual weak primitivity, local 2-transitivity and
geometric $\mathrm{IP}_2$. The first three and $\mathrm{IP}_2$ are not
inferred from an earlier subgroup-intersection check. Residual weak
primitivity quantifies over every nonempty residue, including panels,
and asks for at least one primitive type action of degree at least two.
The exact double-coset maximality criterion permits nonfaithful actions.
Local 2-transitivity checks the actual point-stabilizer orbits in each
panel. Every excluded input below has a recorded failed hypothesis.

\section{The combined sporadic calculation}
Table~\ref{counts} gives one row for each group and rank. Each triple
lists IP, WIP without IP, and failure of WIP at the indicated residue
rank. The last column takes the conjunction over all ranks, not just
the top rank. Every retained class satisfies $\mathrm{IP}_2$.
The machine-readable table additionally retains all 57 literal source
rows, their page and table locators, their original numbers, and the
four precise differences.

\begin{center}\small
\begin{longtable}{llrrrrrr}
\caption{Exact current sporadic roster counts.}\label{counts}\\
\toprule Group&Rank&$N$&$k=3$&$k=4$&$k=5$&$k=6$&Full\\
\midrule\endfirsthead
\toprule Group&Rank&$N$&$k=3$&$k=4$&$k=5$&$k=6$&Full\\
\midrule\endhead
M11&3&13&6/1/6&--&--&--&6/1/6\\
M11&4&11&2/0/9&2/0/9&--&--&2/0/9\\
M12&3&17&7/1/9&--&--&--&7/1/9\\
M12&4&18&4/0/14&3/0/15&--&--&3/0/15\\
M12&5&1&1/0/0&1/0/0&0/0/1&--&0/0/1\\
M22&3&25&5/8/12&--&--&--&5/8/12\\
M22&4&16&5/4/7&3/1/12&--&--&3/1/12\\
M22&5&9&0/1/8&0/1/8&0/1/8&--&0/1/8\\
M23&3&13&5/1/7&--&--&--&5/1/7\\
M23&4&21&5/2/14&3/1/17&--&--&3/1/17\\
M23&5&10&2/1/7&1/1/8&1/0/9&--&1/0/9\\
J1&3&1&0/0/1&--&--&--&0/0/1\\
J1&4&2&1/0/1&1/0/1&--&--&1/0/1\\
J2&3&19&8/3/8&--&--&--&8/3/8\\
J2&4&3&1/0/2&1/0/2&--&--&1/0/2\\
J3&3&19&6/8/5&--&--&--&6/8/5\\
HS&3&30&8/5/17&--&--&--&8/5/17\\
HS&4&41&8/1/32&5/0/36&--&--&5/0/36\\
HS&5&8&1/0/7&1/0/7&0/1/7&--&0/1/7\\
HS&6&1&0/0/1&0/0/1&0/0/1&0/0/1&0/0/1\\
\bottomrule
\end{longtable}

\end{center}
\noindent The entries
for $M_{12}$ rank three and $J_3$ are qualified below. Dashes denote a
residue rank larger than the geometry rank. The $k=6$ entry records the
source's separate sentence about the unique rank-six HS geometry.

The total in this table is 278 representatives, compared with the printed
population 277. This difference does not assert a new global sporadic
classification. It consists of one fewer eligible $M_{12}$ rank-three
class and two additional inequivalent eligible $J_3$ inputs in the
supplied archive. All other numerical source rows agree.

\paragraph{Mathieu inputs.}
The literal $M_{11}$ archive has thirteen rank-three and eleven rank-four
inputs. All twenty-four pass the source geometry hypotheses. Its final
rank-three reorder is retained explicitly. For $M_{12}$ there are twenty
active rank-three, twenty rank-four and one rank-five input. Rank-three
active entries 15,17,20 and rank-four entries 4,7 fail local
2-transitivity; the remaining seventeen, eighteen and one respectively
give the table. All their residues and directional decisions were
computed. The prior complete $M_{12}$ subgroup census and the exact
archive comparison are retained, not rerun; the one missing census
class is treated below.

\paragraph{First and second Janko inputs.}
The three $J_1$ archive families match entire accepted marked families
in the prior complete census: one rank-three and two rank-four cases.
For $J_2$, twenty-five archive entries represent nineteen rank-three
classes and only one of the three rank-four classes. The other two
rank-four cases are supplied by the accepted complete Hall--Janko census,
whose exact representatives are retained. Every archive entry, including
duplicates, has a whole-family correspondence. No new subgroup census
is performed here. Earlier accepted geometry predicates transfer only
across those actual family maps; the shadow decisions are computed on
all twenty-two class representatives.

\paragraph{The two larger Mathieu groups.}
Sixty-six active $M_{22}$ families match fifty retained phase-seven
classes, of ranks $3,4,5$ in numbers $25,16,9$. Six further construction
entries lie outside those sixty-six active inputs; they are separately
accounted for by their failed local 2-transitivity.
Forty-four active $M_{23}$ families match forty-four retained classes,
of ranks $3,4,5$ in numbers $13,21,10$. The old complete census and
residual weak-primitivity conclusions are retained; the maps identify
the actual whole families, and all other current geometry and shadow
checks are performed. Historical publication suffixes are not the
phase-seven class indices.

\paragraph{Higman--Sims.}
All 124 literal active families---44 of rank three, 67 of rank four,
12 of rank five and one of rank six---pass actual source-filter checks.
The exact full-automorphism/type quotient has $30,41,8,1$ classes.
Forty-four classes occur twice and thirty-six once. The full
automorphism group has order $88\,704\,000$. All intersection orders
are necessary rejection invariants only; the actual subgroup
transporters establish positive family identifications. Caching target
stabilizer chains does not alter the permitted transporter group.
The eighty representatives contribute 1,136 base residues including
panels, 836 residues of rank at least two, and 2,204 directional shadow
decisions. These are all used in the table. Completeness beyond the
literal roster remains the classification cited in source reference [21].

\section{The Mathieu exception does not come from a missing input}
The retained complete $M_{12}$ rank-three census has twenty-one classes;
twenty have exact raw-archive correspondents. The remaining census
class 18 is absent from the archive, but it cannot repair the source's
eighteen-geometry locally 2-transitive population. This class number is
our earlier census index, not a historical source name.
\begin{proposition}
Let $G=M_{12}$ and let $(H_1,H_2,H_3)$ be the marked subgroup family
in retained census class 18. Its orders are
\[
 |G|=95040,\qquad (|H_1|,|H_2|,|H_3|)=(24,660,660),
 \qquad |H_1\cap H_2\cap H_3|=6.
\]
The geometry satisfies all of the above eligibility conditions except
local 2-transitivity. It satisfies the full geometric intersection
property. Its panel of type 1 is the action of $A_5$ on the ten
two-element subsets of a five-element set. In particular, it cannot
be added to the locally 2-transitive population.
\end{proposition}
\begin{proof}
The entire marked family is retained in
\path{prior_representatives.g}. The unchanged exact residue procedure
from the common shadow paper checks every base residue, including
panels, and every remaining type. Flag transitivity makes these
representatives exhaustive for residues of the same type. The resulting
certificate verifies firmness, residual connectedness, residual weak
primitivity, and the geometric intersection property in ranks two and
three. These are exact finite group and coset calculations; no property
is inferred solely from the displayed orders.

The failure of local 2-transitivity has a smaller independent proof.
Put $B=H_1\cap H_2\cap H_3$ and $K=H_2\cap H_3$. A bijective
homomorphism $K\longrightarrow A_5$ maps $B$ to
\[
 D=\langle(2\,4\,3),(1\,5)(2\,3)\rangle.
\]
The program \path{identify_panel.g} checks this homomorphism and
that $D$ is exactly the stabilizer of $\{1,5\}$. Thus the right
$B$-cosets in $K$ are identified equivariantly with the ten two-element
subsets. This also gives a proof that is independent of the shadow
procedure: \path{panels.g} constructs those ten right cosets directly
and computes the action of their point stabilizer.

The stabilizer $D$ has three orbits on the two-element subsets. The
fixed subset is $\{1,5\}$. There are three subsets disjoint from it,
and six meeting it in one point. The displayed three-cycle permutes
the complementary three points, while the displayed involution
interchanges 1 and 5 and two of those three points. These generators
act transitively on each of the latter two sets. The subdegrees are
therefore $1,3,6$. Since the point stabilizer is not transitive on
the other nine points, the panel action is not 2-transitive. The
other two panels each have degree two. This proves the stated
failure without relying on a previously recorded eligibility flag.
\end{proof}


Consequently, under the retained census and its exact family
correspondences, the eligible rank-three population remains seventeen,
with counts $7,1,9$. The source prints $7,1,10$ twice, in its rank-three
and full-property tables. The omission and its non-2-transitivity were
already recorded in the earlier census comparison; the complete shadow
decision and the explicit small panel proof make the remaining
obligation precise.

\section{The third Janko roster has two additional classes}
The final active assignment in \texttt{j3.rk3} swaps construction entries
8 and 9. All active indices in the following argument use that final
assignment. The group has order $50\,232\,960$. Every one of the
twenty-one literal inputs passes actual geometry and local-filter tests
on all 147 base residues, including panels. The full automorphism group
has order $100\,465\,920$ and the group has trivial centre. A verified
noninner automorphism therefore completes the identity to a transversal
modulo inner automorphisms. The common simultaneous-conjugacy criterion
tests both possibilities and every type permutation; it assumes no
split extension. Saved positive maps retain the images of every ambient
generator, the actual inner conjugator and the type permutation.
\section{Explicit class separations and intersection results}
For $J\subseteq\{1,2,3\}$ let
$G_J=\bigcap_{j\in J}G_j$, with $G_\varnothing=G$.
The multiset of all eight orders $|G_J|$ is invariant under group
automorphisms and type permutations. On the supplied roster its buckets
are the singleton entries and
\[
 \{7,8\},\quad \{9,10,11\},\quad
 \{13,14\},\quad \{16,17\}.
\]
There are $16$ buckets. Actual complete-family transporters identify
$13$ with $14$, and $16$ with $17$; their type permutations,
automorphism images of every ambient generator, and the conjugating
permutations are retained as exact data. The other repeated buckets
split, with the following independent witnesses.

Entries $7$ and $8$ each have a unique parabolic of order $96$; these
parabolics contain respectively $27$ and $15$ involutions. Entry $9$ has
two parabolics of order $96$, each with $27$ involutions, whereas entries
$10$ and $11$ have two such parabolics each with $15$. These actual counts
are obtained from complete conjugacy-class rosters in the small
parabolic groups; group isomorphisms preserve them. Finally entry $10$
satisfies full IP, while entry $11$ fails WIP. The complete shadow
certificates, with the failed directional intersections, therefore
separate this last pair as well. These tests independently support the
negative transporter decisions and give $16+1+2=19$ distinct supplied
classes without relying on historical names.

\begin{proposition}
Modulo entire-family automorphisms of the retained $J_3$ action and type
permutations, the supplied roster has $19$ eligible classes. Every class
satisfies $\mathrm{IP}_2$. Full IP, equivalently $\mathrm{IP}_3$ here,
holds in six classes; eight more satisfy WIP; five fail WIP.
\end{proposition}
\begin{proof}
The order profiles separate different buckets. The explicit transporters
identify exactly the two stated pairs, and the preceding invariants
separate the other entries, so the quotient has $19$ classes. All supplied
entries were checked for the source geometry hypotheses, so these
predicates transfer across the actual family isomorphisms.

For each representative, the accepted residue theorem enumerates all
seven proper fixed-type subsets. Flag transitivity transports each base
residue to all flags of those types. Panels are checked for geometry
properties and only ranks two and three enter IP and WIP. This gives
$133$ base residues including panels, $76$ of rank at least two, and
$171$ directional shadow decisions. Their actual all/exists quantifiers
are reconstructed from the recorded values. The positive active
representatives are
\[
 2,3,4,7,8,10;
\]
the weak-only representatives are
\[
 1,5,6,12,16,18,19,21;
\]
and the nonweak representatives are
\[
 9,11,13,15,20.
\]
All rank-two residues satisfy IP, so the rank-three decision and full
residual decision agree. This proves the asserted supplied-roster counts.
\end{proof}


Thus the supplied-roster split is $6,8,5$, whereas the source prints
$5,8,4$. The global earlier classification in source reference [19]
is retained separately; no unprovided input or historical suffix is
identified by the number nineteen.

\section{Historical labels and the corrected source contract}
The complete companion \emph{Historical names and exact shadow decisions
for sporadic geometries} supplies the source-profile argument, the bounded
correlation and residue computations, and every compatible family for
229 of the 232 positive occurrences in Tables 2,4,6,8. It uses retained
complete populations and invariant data read from the primary diagram
supplements; the desired shadow answer is never used to choose a model.
The README's HS4.1 pointer is corrected to the uniquely matching actual
profile, rather than assigning its claimed answer to the wrong input.

The three occurrences called $J_1$4.1 are stated mathematically through
the complete two-class rank-four population. The class with type sizes
$\{266,1463,1463,2926\}$ and Borel order2 has full IP. The other class,
with type sizes $\{266,266,1463,2926\}$ and Borel order2, fails WIP
in ranks3 and4. These are exact results of the accepted complete
census and whole-family certificates. Their literal identification with
the historical ordinal4.1 in the 1998 chapter remains unverified.
It is retained as historical metadata, not silently asserted or used as
a mathematical premise. All three corresponding aggregate counts are
recovered. Thus the full corrected principal contract is explicit,
without claiming equality of the literal original and replacement
labelled theorem sets.

\section{The Suzuki application}
The following local enumeration uses the whole thin-triple and residue
theory in the common paper. It retains the specified degree-65 Suzuki
model and the earlier classification assertion that the relevant
rank-three RWPRI geometries are thin; it does not claim recovery of
the unavailable earlier classification article.
\section*{A finite enumeration with an explicit orbit reduction}
The ambient permutation group is the previously validated degree-65 model
used in the rank-two Suzuki calculation. We retain that model and GAP's
exact finite group algorithms. The computation checks order 29120 and
simplicity, one conjugacy class containing all 455 involutions, and an
automorphism group of order 87360. Its action on those involutions has
trivial kernel. It also computes representatives of all maximal subgroups,
whose orders are $448,52,20,14$. These are explicit computational premises;
their implementations have not been formalized in Lean here.

Let $\mathcal A$ be the computed automorphism group acting on the 455
involutions, numbered in a fixed order. Fix the first involution. For its
stabilizer $\mathcal A_1$, take one representative $b$ from each orbit on
the other involutions. For each $b$, take one $c$ from each orbit of
$\mathcal A_{1,b}$ on the remaining involutions. This gives all ordered
triples up to $\mathcal A$, exactly once.

To forget the type order, normalize each of the six permutations of a
triple by the same stabilizer procedure and retain the smallest resulting
pair of indices $(b,c)$. The normalization is independent of transporter
choices: after placing the first entry at 1, any two choices differ by
$\mathcal A_1$; after placing the second entry at $b$, they differ by
$\mathcal A_{1,b}$. Thus the final orbit of the third entry is fixed.
Equality of the minimum pairs is equivalent to equivalence under
$\mathcal A$ and a permutation of the three types.

For each retained triple the program checks generation, the exact
subgroup intersections, all three product equalities~\eqref{U-products},
and the primitivity conditions of Lemma~\ref{U-thin-recognize}. Top primitivity
is checked by conjugacy to the computed maximal subgroups, not inferred
from a subgroup order. The resulting counts are
\begin{center}\begin{tabular}{lr}\toprule
Stage & Count\\\midrule
Ordered triples modulo $\mathcal A$ & 1083\\
After forgetting the three type names & 195\\
Generating triples & 186\\
Triples satisfying all marked subgroup intersections & 186\\
Flag-transitive thin geometries & 183\\
Residually weakly primitive geometries & 147\\\bottomrule
\end{tabular}\end{center}
The three rejected generating intersection systems show why a subgroup
intersection test alone does not establish the required geometry.

For these thin connected geometries the type-preserving automorphism
group is exactly $G$. Indeed an automorphism fixing one chamber fixes its
unique neighbors across every panel, and then every chamber by connectedness.
The full type-preserving group is consequently semiregular on chambers,
while $G$ is already transitive. It can have no more elements than $G$.
Thus an isomorphism between two such geometries induces an automorphism
of $G$ and a permutation of types, exactly the equivalence used above.

\section*{Shadow tests and their limits}
For each retained triple, the common shadow theorem expresses every
shadow as an orbit of a chamber stabilizer intersection in $H_i\backslash G$.
It checks all normalized object--flag pairs and every allowable witness
flag. The implementation is the earlier reviewed shadow implementation,
with the ambient order and input triples changed. It caches conjugate
coset actions: if $H'=H^g$, the base point changes from $p$ to $p^g$.
The action is unchanged, and the stabilizer of the new point is checked
to equal $H'$. No shadow or witness is removed by this optimization.

The full eligible-extension test is repeated in all three coset actions
for each geometry, providing 1764 flag-extension identities in addition
to the product tests above. Each of the 441 type tests is recorded.
The resulting numbers are 27 with all three types passing, 48 with some
but not all types passing, and 72 with no type passing. All rank-two
residues are polygons of lengths $2m_i$ with $m_i\in\{2,5,7,13\}$.
For $m_i>2$ there is no quadrilateral: shadows of distinct opposite-type
vertices meet in at most one object, which itself is a witness whenever
present. If $m_i=2$, the residue is the complete bipartite graph $K_{2,2}$;
the two opposite-type shadows are the entire target type, with the empty
flag as witness. Flags already containing the target type have singleton
shadows and use that object as witness; the empty input flag has the
whole type as shadow. This proves the rank-two intersection property in
both cases. Rank-one residues are automatic.
Consequently the rank-three answers are also the full residual answers.

The orbit and product arguments enumerate all \emph{thin} geometries
with the stated hypotheses. They do not prove that every residually
weakly primitive rank-three geometry of $\operatorname{Sz}(8)$ is thin.
That separate classification assertion is retained from the earlier
classification cited by Jacobs and Leemans. We claim no recovery of its
unavailable manuscript. The present result is the reported shadow split
on its 147 thin inputs.


The five eligible square-residue cases are included, not discarded as
degenerate polygons. Among all 183 thin flag-transitive cases there are
seven with a square residue: five pass RWPRI and two fail it. The seven
is compatible with the previously accepted Topic D regular Suzuki
polyhedron count $P(3)=7$ up to duality. That older count and its proof
remain attributed to that earlier work; they are a comparison, not a
new premise or new result obtained from matching this number.

\section{Transfer to inductively minimal geometries}
We use the precise convention recalled in [BC98]: a pair consisting of
a finite firm rank-$r$ geometry with connected diagram and a
flag-transitive group $G$ is called \emph{minimal} when
$|G|\leq(r+1)!$. The word refers to this numerical threshold, not to
minimality in a subgroup ordering. Inductive minimality requires that
for each connected type subgraph $J$ and every flag $F$ of complementary
type, the residue pair $(\Gamma_F,G_F)$ is minimal in this sense.
We retain the foundational classification of Buekenhout, Cara and Dehon,
in the precise form recalled by Buekenhout and Cara~[BC98]: an inductively
minimal pair of positive rank $r$ has group $\Sym(r+1)$; its connected
diagram has no induced cycle longer than three, each edge belongs to a
unique maximal clique, and each vertex belongs to one or two maximal
cliques. Every such connected diagram determines one pair up to
isomorphism. The one-vertex diagram is treated separately. The proof of
this foundational classification is not replaced here.

These diagrams are precisely line graphs of trees. To see the required
direction, form a vertex for each maximal clique. A diagram vertex in
two cliques gives an edge between their clique vertices; a vertex in
one clique gives an edge to a new leaf. Distinct cliques meet in at most
one vertex, because a common edge would violate uniqueness of its
clique. The constructed graph is simple and connected, and its line
graph is the original diagram. A shortest cycle of length at least four
would give an induced cycle on its edge labels in the diagram. A triangle
would place three pairwise adjacent labels in one clique, contradicting
unique ownership of the pairwise edges by the three cliques already
represented at the cycle vertices. The constructed graph is therefore
a tree. Its reconstruction is inverse to taking the nontrivial vertex
stars of a tree as line-graph cliques. The one-vertex diagram corresponds
to the one-edge tree.

For the geometry in Lemma~\ref{U-forest}, two remaining edges give either
a six-cycle residue or $K_{2,2}$, according as they meet or are disjoint.
Its diagram is consequently the line graph of $T$. Every residue with
connected diagram and $m$ remaining types comes from a single subtree
on $m+1$ vertices, and its group is $\Sym(m+1)$. Its order is exactly
$(m+1)!$, satisfying the stated source definition of minimality.
Thus $\Gamma(T)$ is inductively minimal.
The retained uniqueness identifies every inductively minimal pair with
one of these tree geometries.

\begin{corollary}
An inductively minimal geometry has the intersection property exactly
when its diagram is linear. In particular this holds for the
$\Sym(n)$ geometries of rank $n-1$ with $2\leq n\leq7$.
\end{corollary}
\begin{proof}
Apply the retained classification, the construction just proved, and
Theorem~\ref{U-path}. The line graph of a tree is a path exactly when the
tree itself is a path: a vertex of degree at least three would give a
triangle in the line graph.
\end{proof}



\section{Exact executable, retained and formal boundaries}
The complete finite programs and actual input data are mathematical
dependencies for the computational conclusions. Their role inventory
distinguishes source parsing, actual geometry eligibility, family
transport, shadow decisions and result aggregation from optional
validation and operational rendering. Successful log summaries alone
are not substituted for the full executable files. The old accepted
subgroup-census programs, models, completeness arguments and scoped
review records remain dependencies wherever the preceding applications
retain them; those computations are not silently free prior results.
GAP's exact finite-group algorithms and installed catalogue inputs are
trusted computational infrastructure, with their pinned versions and
limits retained.

No unchanged large census or completed full eligibility batch needs to
be rerun to assemble this manuscript. The current aggregation recomputes
every all/exists quantifier and all 57 source comparisons directly from
the accepted directional records. For $J_3$, the optional five-program
route recomputes family transport, saved-map data, small local invariants
and shadows; it retains the four actual eligibility batches. A stale
opening docstring in the frozen runner says maximality is rerun; its
README and mathematical report give the correct boundary, retained here.

The source prints concrete Magma programs in Section 6 and an executable
algorithm in Figure 1. Future manuscript-only lengths must apply the
same exclusion rule to these listings and to ours, while retaining
mathematical explanation on both sides. Whole required software and
data remain visible separately; missing source software costs are
unknown, never zero. Table 9 and the Suzuki twelve-second statement
are historical execution reports, not new mathematical or performance
claims. This assembly supplies no length ratio.

The accepted four-lemma covering formal module certifies its exact
set/flag covering equivalence. The inherited Boolean-shadow module
proves the interval description and at-most-two-element witness for
finite query families. Neither proves the whole orbit algorithm,
sporadic computations, thin Suzuki census, tree identification or the
retained classifications in Lean. Their complete files and exact
declaration scope remain separately bound. There is no new conjecture
credit: the all-degree tree result was already attributed by the source
to Cara and Leemans.

\begin{thebibliography}{9}
\bibitem[JL]{jl} P.~Jacobs and D.~Leemans, An algorithmic analysis of the
intersection property, \emph{LMS J. Comput. Math.} 7 (2004), 284--299.
\bibitem[BC98]{bc} F.~Buekenhout and P.~Cara, Some properties of
inductively minimal geometries, \emph{Bull. Belg. Math. Soc. Simon Stevin}
5 (1998), 213--219. The foundational 1997 classification is retained
in the precise form recalled there.
\bibitem[CL05]{cl} P.~Cara and D.~Leemans, On inductively minimal
geometries that satisfy the intersection property,
\emph{J. Combin. Theory A} 111 (2005), 327--330.
The source attributes the uniform conclusion to this prior work;
the unavailable full article receives no original credit here.
\end{thebibliography}
\end{document}
